\section{Evidencia temprana: Machine Translation, Parsing y Attention Pattern}

Que un mecanismo sea razonable no significa que el modelo sea eficaz. El Transformer original mostró primero su calidad y su eficiencia de entrenamiento en WMT machine translation, y después se trasladó a constituency parsing; la attention visualization, por su parte, ofrece indicios locales de interpretabilidad. Esta sección separa tres tipos de evidencia: benchmark improvement, cross-task transfer y qualitative pattern.

\subsection{Machine Translation: mejoran a la vez la calidad y la eficiencia de entrenamiento}

El artículo original compara BLEU, costo de entrenamiento y tamaño del modelo en WMT14 English--German y English--French. Al leer la tabla, primero hay que confirmar si se trata de single model o ensemble y cuál es el presupuesto de entrenamiento, y solo después comparar la calidad. Lo importante del Transformer no es que gane unas cuantas décimas de BLEU, sino que alcanza o supera a los sistemas de la época con una arquitectura más fácil de paralelizar.

\lecturefigure{slide-24-sota-machine-translation.jpg}{Resultados de machine translation del Transformer original}{Diapositiva V024 recuperada de la grabación oficial; 00:28:54}

\begin{knowledgebox}{Lectura de la figura: primero alinear el tipo de modelo y el costo de entrenamiento}
La tabla incluye a la vez baselines recurrent, convolutional, single-model y ensemble. Primero se comparan configuraciones del mismo tipo y después se revisan los FLOPs o el tiempo de entrenamiento; la evidencia central es que el Transformer obtiene un BLEU alto con un costo de entrenamiento menor. BLEU sigue siendo una corpus-level overlap metric: no abarca toda la semántica, la robustez ni la calidad en un despliegue real.
\end{knowledgebox}

\subsection{Generalization to Parsing: la arquitectura no sirve solo para traducción}

En constituency parsing, la estructura de entrada y salida es distinta de la de translation. Aplicar la misma architecture a parsing pone a prueba si el motif de cómputo puede reutilizarse entre tareas, en lugar de limitarse a memorizar rasgos de la traducción. Los resultados respaldan que la pila self-attention + FFN es un backbone general de sequence transduction.

\lecturefigure{slide-25-generalization-to-parsing.jpg}{El Transformer generaliza de machine translation a constituency parsing}{Diapositiva V025 recuperada de la grabación oficial; 00:29:45}

\begin{knowledgebox}{Lectura de la figura: el resultado entre tareas valida la architecture reuse}
Primero se revisan los parsing benchmarks, como WSJ, y sus baselines; después se distingue entre «el mismo block puede entrenarse» y «no se necesitan task-specific data». Parsing conserva su propia representation de salida, sus datos y su objective; la figura respalda la backbone consolidation, pero eso no equivale a que todas las tareas solo requieran un prompt zero-shot.
\end{knowledgebox}

\subsection{Attention Map: un Pattern legible no es una explicación completa}

Algunas heads muestran dependencias de larga distancia, patrones sintácticos o de copia, por lo que se vuelven una ventana para que los investigadores entiendan el modelo. Sin embargo, el attention weight solo describe los coeficientes de enrutamiento de la layer/head actual: no incluye el value content, la residual contribution ni la nonlinear transformation posterior, y tampoco demuestra directamente importancia causal.

\lecturefigure{slide-26-interpretable-attention-patterns.jpg}{Patterns de larga distancia y sintácticos visibles en un attention map}{Diapositiva V026 recuperada de la grabación oficial; 00:30:48}

\begin{warningbox}{«Parece gramática» no significa que el modelo dependa causalmente de esa Head}
Para leer la figura, primero se confirman los tokens de los ejes horizontal y vertical, la escala de color y la head/layer; después se observa si los pesos altos son estables entre ejemplos. Para demostrar un mecanismo se requieren además ablation, activation patching, counterfactual intervention o value-path analysis; un solo heatmap solo puede aportar una hypothesis.
\end{warningbox}

\teachervoice{El ponente mantiene una postura mesurada respecto a la attention interpretability: algunos patterns son muy vistosos y en efecto parecen long-distance dependency, pero la capacidad de la architecture puede deberse a la acción conjunta del enrutamiento de información y de la enorme FFN. No conviene elevar una sola figura a «proceso de pensamiento del modelo».}

\subsection{Resumen de la sección}

La early evidence se divide en tres niveles: la tabla de translation demuestra calidad y eficiencia, parsing demuestra que el backbone es transferible y el attention map ofrece patterns internos que pueden investigarse. Los tres niveles son importantes, pero respaldan conclusiones de distinta solidez; un apunte de calidad debe conservar esta evidence discipline.
